% GENERATED FILE. Edit canonical lessons, then run npm run generate:books.
% canonical-source-hash: fnv1a64:6e219bd8093f5ad0
% canonical-lessons: AR-C537-takhyim, AR-C537-ghaws, AR-C537-taswir, AR-C537-fakhkhar, AR-C537-kalimat-mutaqatia

\chapter{Hobbies indoors and out}
\label{ch:ar-hobbies-indoors-and-out}
\hlchaptermodality{537}

\begin{chapteropening}
\textbf{By the end of this chapter:} \emph{I can talk about camping, diving, photography, pottery and a crossword.}

Complete the last lesson of chapter 537: I can talk about camping, diving, photography, pottery and a crossword.
\end{chapteropening}

\section[\texorpdfstring{\ar{تخييم} (takhyīm)}{تخييم (takhyīm)}]{\ar{تخييم} (takhyīm) — camping}

\label{lesson:AR-C537-takhyim}



\begin{quote}
Before the new one: say the Arabic for to unwind, then the Arabic for to enjoy.
\end{quote}

\subsection*{You'll want to know: \ar{تخييم}}
\textbf{\ar{تخييم}} — \emph{takhyīm} — \textquotedblleft{}camping\textquotedblright{}.

\begin{quote}
\ar{نحب} \ar{التخييم} \ar{في} \ar{الصحراء}.
\end{quote}

\begin{quote}
\emph{nuḥibbu at-takhyīm fī aṣ-ṣaḥrāʾ.}
\end{quote}

\begin{quote}
We love camping in the desert.
\end{quote}

\begin{cousinweb}
\textbf{\ar{تخييم}} (\emph{takhyīm}) is built on the native Arabic root \emph{kh-y-m}, \textquotedblleft{}a tent\textquotedblright{} (Lane). Khayma, a tent, which you know, is from the same root.
\end{cousinweb}

\begin{grammarlens}[title={What you've built}]
One more word for describing what you lived and felt.
\end{grammarlens}

\subsection*{Your turn}
\begin{itemize}
\raggedright
  \item \emph{Say it:} \emph{takhyīm}
  \item \emph{Say it:} \emph{takhyīm}, once more
  \item \emph{Say it:} \emph{nuḥibbu at-takhyīm fī aṣ-ṣaḥrāʾ.} — We love camping in the desert.
  \item \emph{Recall:} say the Arabic for to unwind, then the Arabic for to enjoy, then say \emph{takhyīm} again
\end{itemize}

\begin{tcolorbox}[breakable,colback=teal!4,colframe=teal!35!black,title={Before you move on}]
What does \ar{تخييم} mean? (\textquotedblleft{}Camping\textquotedblright{}.) Say it once more.
\end{tcolorbox}
\section[\texorpdfstring{\ar{غوص} (ghawṣ)}{غوص (ghawṣ)}]{\ar{غوص} (ghawṣ) — diving}

\label{lesson:AR-C537-ghaws}



\begin{quote}
Before the new one: say the Arabic for to enjoy, then the Arabic for camping.
\end{quote}

\subsection*{You'll want to know: \ar{غوص}}
\textbf{\ar{غوص}} — \emph{ghawṣ} — \textquotedblleft{}diving\textquotedblright{}.

\begin{quote}
\ar{تعلمت} \ar{الغوص} \ar{في} \ar{الصيف} \ar{الماضي}.
\end{quote}

\begin{quote}
\emph{taʿallamtu al-ghawṣ fī aṣ-ṣayf al-māḍī.}
\end{quote}

\begin{quote}
I learned diving last summer.
\end{quote}

\begin{cousinweb}
\textbf{\ar{غوص}} (\emph{ghawṣ}) is built on the native Arabic root \emph{gh-w-ṣ}, \textquotedblleft{}to dive, to plunge\textquotedblright{} (Lane). Ghawṣ is the plunge under water.
\end{cousinweb}

\begin{grammarlens}[title={What you've built}]
One more word for describing what you lived and felt.
\end{grammarlens}

\subsection*{Your turn}
\begin{itemize}
\raggedright
  \item \emph{Say it:} \emph{ghawṣ}
  \item \emph{Say it:} \emph{ghawṣ}, once more
  \item \emph{Say it:} \emph{taʿallamtu al-ghawṣ fī aṣ-ṣayf al-māḍī.} — I learned diving last summer.
  \item \emph{Recall:} say the Arabic for to enjoy, then the Arabic for camping, then say \emph{ghawṣ} again
\end{itemize}

\begin{tcolorbox}[breakable,colback=teal!4,colframe=teal!35!black,title={Before you move on}]
What does \ar{غوص} mean? (\textquotedblleft{}Diving\textquotedblright{}.) Say it once more.
\end{tcolorbox}
\section[\texorpdfstring{\ar{تصوير} (taṣwīr)}{تصوير (taṣwīr)}]{\ar{تصوير} (taṣwīr) — photography}

\label{lesson:AR-C537-taswir}



\begin{quote}
Before the new one: say the Arabic for camping, then the Arabic for diving.
\end{quote}

\subsection*{You'll want to know: \ar{تصوير}}
\textbf{\ar{تصوير}} — \emph{taṣwīr} — \textquotedblleft{}photography\textquotedblright{}.

\begin{quote}
\ar{هواية} \ar{أختي} \ar{التصوير}.
\end{quote}

\begin{quote}
\emph{hiwāyat ukhtī at-taṣwīr.}
\end{quote}

\begin{quote}
My sister's hobby is photography.
\end{quote}

\begin{cousinweb}
\textbf{\ar{تصوير}} (\emph{taṣwīr}) is built on the native Arabic root \emph{ṣ-w-r}, \textquotedblleft{}to form, to shape\textquotedblright{} (Lane). Ṣūra, a picture, and ṣawwara, to photograph, which you know, share the root.
\end{cousinweb}

\begin{grammarlens}[title={What you've built}]
One more word for describing what you lived and felt.
\end{grammarlens}

\subsection*{Your turn}
\begin{itemize}
\raggedright
  \item \emph{Say it:} \emph{taṣwīr}
  \item \emph{Say it:} \emph{taṣwīr}, once more
  \item \emph{Say it:} \emph{hiwāyat ukhtī at-taṣwīr.} — My sister's hobby is photography.
  \item \emph{Recall:} say the Arabic for camping, then the Arabic for diving, then say \emph{taṣwīr} again
\end{itemize}

\begin{tcolorbox}[breakable,colback=teal!4,colframe=teal!35!black,title={Before you move on}]
What does \ar{تصوير} mean? (\textquotedblleft{}Photography\textquotedblright{}.) Say it once more.
\end{tcolorbox}
\section[\texorpdfstring{\ar{فخار} (fakhkhār)}{فخار (fakhkhār)}]{\ar{فخار} (fakhkhār) — pottery}

\label{lesson:AR-C537-fakhkhar}



\begin{quote}
Before the new one: say the Arabic for diving, then the Arabic for photography.
\end{quote}

\subsection*{You'll want to know: \ar{فخار}}
\textbf{\ar{فخار}} — \emph{fakhkhār} — \textquotedblleft{}pottery\textquotedblright{}.

\begin{quote}
\ar{صنعت} \ar{أختي} \ar{إناء} \ar{من} \ar{الفخار}.
\end{quote}

\begin{quote}
\emph{ṣanaʿat ukhtī ināʾ min al-fakhkhār.}
\end{quote}

\begin{quote}
My sister made a pot out of clay.
\end{quote}

\begin{cousinweb}
\textbf{\ar{فخار}} (\emph{fakhkhār}) is built on the native Arabic root \emph{f-kh-r}, \textquotedblleft{}to boast, to glory in\textquotedblright{} (Lane). Lane gives fakhkhār, under the same root as fakhūr, proud, as baked pottery and baked clay.
\end{cousinweb}

\begin{grammarlens}[title={What you've built}]
One more word for describing what you lived and felt.
\end{grammarlens}

\subsection*{Your turn}
\begin{itemize}
\raggedright
  \item \emph{Say it:} \emph{fakhkhār}
  \item \emph{Say it:} \emph{fakhkhār}, once more
  \item \emph{Say it:} \emph{ṣanaʿat ukhtī ināʾ min al-fakhkhār.} — My sister made a pot out of clay.
  \item \emph{Recall:} say the Arabic for diving, then the Arabic for photography, then say \emph{fakhkhār} again
\end{itemize}

\begin{tcolorbox}[breakable,colback=teal!4,colframe=teal!35!black,title={Before you move on}]
What does \ar{فخار} mean? (\textquotedblleft{}Pottery\textquotedblright{}.) Say it once more.
\end{tcolorbox}
\section[\texorpdfstring{\ar{كلمات} \ar{متقاطعة} (kalimāt mutaqāṭiʿa)}{كلمات متقاطعة (kalimāt mutaqāṭiʿa)}]{\ar{كلمات} \ar{متقاطعة} (kalimāt mutaqāṭiʿa) — a crossword}

\label{lesson:AR-C537-kalimat-mutaqatia}



\begin{quote}
Before the new one: say the Arabic for photography, then the Arabic for pottery.
\end{quote}

\subsection*{You'll want to know: \ar{كلمات} \ar{متقاطعة}}
\textbf{\ar{كلمات} \ar{متقاطعة}} — \emph{kalimāt mutaqāṭiʿa} — \textquotedblleft{}a crossword\textquotedblright{}.

\begin{quote}
\ar{تحل} \ar{أمي} \ar{الكلمات} \ar{المتقاطعة} \ar{كل} \ar{صباح}.
\end{quote}

\begin{quote}
\emph{taḥullu ummī al-kalimāt al-mutaqāṭiʿa kull ṣabāḥ.}
\end{quote}

\begin{quote}
My mother does the crossword every morning.
\end{quote}

\begin{cousinweb}
\textbf{\ar{كلمات} \ar{متقاطعة}} (\emph{kalimāt mutaqāṭiʿa}), \textquotedblleft{}crossing words\textquotedblright{}, is a modern calque built from \emph{kalima}, a word, and \emph{mutaqāṭiʿa}, \textquotedblleft{}cutting across one another\textquotedblright{}, from the native root \emph{q-ṭ-ʿ}, \textquotedblleft{}to cut\textquotedblright{} (Lane).
\end{cousinweb}

\begin{grammarlens}[title={What you've built}]
One more word for describing what you lived and felt.
\end{grammarlens}

\subsection*{Your turn}
\begin{itemize}
\raggedright
  \item \emph{Say it:} \emph{kalimāt mutaqāṭiʿa}
  \item \emph{Say it:} \emph{kalimāt mutaqāṭiʿa}, once more
  \item \emph{Say it:} \emph{taḥullu ummī al-kalimāt al-mutaqāṭiʿa kull ṣabāḥ.} — My mother does the crossword every morning.
  \item \emph{Recall:} say the Arabic for photography, then the Arabic for pottery, then say \emph{kalimāt mutaqāṭiʿa} again
\end{itemize}

\begin{tcolorbox}[breakable,colback=teal!4,colframe=teal!35!black,title={Before you move on}]
What does \ar{كلمات} \ar{متقاطعة} mean? (\textquotedblleft{}A crossword\textquotedblright{}.) Say it once more.
\end{tcolorbox}
